\uuid{ya2y}
\titre{Variables continues}
\auteur{worms}
\organisation{exo7}
\licence{CC BY-NC-SA 4.0}
\datecreate{2017-08-28}
\niveau{L2}
\module{Probabilité et statistique}
\chapitre{Probabilité continue}
\sousChapitre{Densité de probabilité}
\difficulte{}
\theme{Variables continues}
\exo7qcmid{worms.worms.040}

\contenu{
\begin{question}
\qclassification{261.01}
Si on considère $X$ une variable aléatoire dont la densité est la fonction représentée ci-dessous, laquelle/lesquelles de ces affirmations est/sont vraie(s) ?

%[[image]]
\qimage{images/png/ya2y-1.png}

%\begin{center}
%\centerline{\includegraphics[width=7.cm,height=4.cm]{GraphepourQCM3.jpg}}
%\end{center}

\begin{answers}
\bad{La médiane de $X$ est positive.}
\good{$\mathbb{P}(X<2) > \mathbb{P}(X>-1)$}
\bad{$\mathbb{P}( \, \ln(X+3) >0 \, ) < 1/2$}
\good{$\mathbb{P}(X>x) < 1/2$ pour tout $x>0$.}
\end{answers}
\begin{explanations}
Pour l'affirmation $\mathbb{P}( \, \ln(X+3) >0 \, ) < 1/2$, on peut répondre non (sans calculs) car $X$ est toujours $\geq -2$, donc $\ln(X+3)$ est toujours $\geq \ln(1)=0$.
\end{explanations}
\end{question}
}
